\textbf{结论名称：} 垂直等长与 $90^\circ$ 旋转坐标关系

\textbf{英文名称：} perpendicular\_equal\_segments\_rotation\_coordinates

\textbf{适用条件：} 平面上三点 $B,P,Q$ 满足 $BP\perp BQ$ 且 $BP=BQ>0$，两条线段有公共端点 $B$。

\textbf{变量说明：} $B(a,b),\ P(x_1,y_1),\ Q(x_2,y_2)$；先令 $u=x_1-a,\ v=y_1-b$，即从 $B$ 出发到 $P$，水平走 $u$、竖直走 $v$。

\textbf{核心公式：}
\[
\boxed{\begin{aligned}
\text{逆时针 }90^\circ:\ &(x_2-a,\ y_2-b)=(-v,\ u),\\
\text{顺时针 }90^\circ:\ &(x_2-a,\ y_2-b)=(v,\ -u).
\end{aligned}}
\]
特别地，$\boxed{|x_2-a|=|y_1-b|}$，即旋转后的水平变化量，其绝对值等于旋转前的竖直变化量。

\textbf{识别信号：} 同一顶点发出两条线段，\textbf{垂直而且等长}；或图中出现以公共点为直角顶点的等腰直角三角形。不要将“仅仅垂直”误当作足够条件。

\textbf{使用场景：} 点绕定点转 $90^\circ$、等腰直角三角形与正方形、解析几何中垂直等长的坐标消元、椭圆等曲线上的点与定直线上的点相配合。

\textbf{简单例子：} $B(5,0)$，$P(x_1,y_1)$，$Q(6,y_2)$，且 $BP\perp BQ$、$BP=BQ$。因为 $x_2-a=6-5=1$，所以 $|y_1|=1$，无需列斜率、距离方程。注意还需根据旋转方向确定 $y_2$。
